\section{File \texttt{input.postproc}}
\subsection{Purpose}
\texttt{input.postproc} is a text file specifying how the \swPostProcess executable is to be run for extracting results. 

\subsection{Format description}
\texttt{<input.postproc>} is a text file. It contains one keyword per line, with format
\begin{lstlisting}
<KEYWORD> <VALUE>
\end{lstlisting}
The order of the keywords does not matter. Empty lines or lines starting with hash '\#' or equal '=' signs are considered as comment lines.\\
Note that
\begin{itemize}
\item The keywords and its value must be separated by at least one space or tabulation characters.
\item The spelling of the keywords is not case sensitive.\\
\item If a given keyword occurs several times in the file, the last read value overwrites the previously specified ones.
\item All physical parameters (e.g. gravitational acceleration, etc.) are in S.I. units.
\end{itemize}

\subsection{List of allowed keywords}

\begin{description}

\item[compute\_max] \defaultvalue{0} Set if the maximal values have to be extracted from the existing results. Allowed values:
\begin{description}
\item[0] no computation of maximum values.
\item[1] maximum values are computed.
\end{description}
The maximum values and corresponding times are stored in files \texttt{output/ResultHydro\_Max.txt} and \texttt{output/ResultHydro\_DateOfMax.txt} respectively.

\item[first\_submersion] \defaultvalue{0} Set if the dates of first submersion have to be extracted from the existing results. Allowed values:
\begin{description}
\item[0] no computation of dates of first submersion.
\item[1] date of first submersion are computed.
\end{description}
The dates of first submersion are stored in file \texttt{output/Submersion\_Durations.txt}.

\item[h\_min] \defaultvalue{0.01} Set the minimum water depth at which a cell is considered flooded for the computation of submersion durations. Allowed values:
\begin{description}
\item Any number greater than or equal to zero
\end{description}

\item[input\_format] \defaultvalue{txt} Set the format of result file(s) to be used to compute the maximal values and submersion values. Allowed values:
\begin{description}
\item[bin] Binary file \texttt{ResultHydro.bin} is used to read the simulation results.
\item[txt] Text files \texttt{ResultHydro**.txt} and/or \texttt{ResultTime**.txt} are used to read the simulation results.
\end{description}

\item[input\_max] \defaultvalue{all} Set the files to be used to compute the maximal values. Allowed values:
\begin{description}
\item[resulthydro] maximum values are computed using the map files \texttt{ResultHydro**.txt} alone. These are the maps stored periodically, with time intervals specified by the \texttt{dt\_map} keyword in the \texttt{input.solv} file.
\item[resulttime] maximum values are computed using the extra storage maps \texttt{ResultTime**.txt} files only. The extra storage times are defined in the \texttt{extraction\_times\_list.txt} file. 
\item[all] maximum values are computed using both regular and extra storage maps \\\texttt{ResultHydro**.txt} and \texttt{ResultTime**.txt}.
\end{description}
NB: \texttt{ResultTime**.txt} file names are expected to comply with the following format \\\texttt{ResultTime\_p????d??h??m??s???\_<custom\_string>.txt} where \texttt{<custom\_string>} is set by the user.

\item[input\_submersion] \defaultvalue{all} Set the files to be used to compute the submersion durations and date of first submersion. Allowed values:
\begin{description}
\item[resulthydro] submersion values are computed using the map files \texttt{ResultHydro**.txt} alone. These are the maps stored periodically, with time intervals specified by the \texttt{dt\_map} keyword in the \texttt{input.solv} file.
\item[resulttime] submersion values are computed using the extra storage maps \texttt{ResultTime**.txt} files only. The extra storage times are defined in the \texttt{extraction\_times\_list.txt} file. 
\item[all] submersion values are computed using both regular and extra storage maps \\\texttt{ResultHydro**.txt} and \texttt{ResultTime**.txt}.
\end{description}
NB: \texttt{ResultTime**.txt} file names are expected to comply with the following format \\\texttt{ResultTime\_p????d??h??m??s???\_<custom\_string>.txt} where \texttt{<custom\_string>} is set by the user.

\item[submersion\_durations] \defaultvalue{none} Set which submersion durations have to be computed. Allowed values:
\begin{description}
\item[none] no computaion of submersion durations.
\item[max] maximum submersion durations are computed: it's the maximum uninterrupted time between a submersion and a drying.
\item[total] total submersion durations are computed: it's the time between the first submersion and the last drying during the simulation.
\item[all] both maximum and total submersion durations are computed.
\end{description}
The submersion durations are stored in file \texttt{output/Submersion\_Durations.txt}.

\end{description}

\subsection{Example}
\begin{verbatim}
= unread line
# sign to comment a full line : (# =)
=======
compute_max   1
input_max     all
input_format  txt
========
\end{verbatim}
